\origpage{80}{ROBERTSON AND SEYMOUR}
\noindent let $P_1$, $P_2$ be the two paths of $C$ between $a$ and $b$. Now $P_1 \cup P_2$, $P_2 \cup P$, $P \cup P_1$ are all circuits, and all have length $< 2\theta$, since $\Gamma'$ is confined. But $U(P) \not\subseteq \operatorname{ins}(P_1 \cup P_2) = \operatorname{ins}(C)$, and so by the second slope axiom, one of $\operatorname{ins}(P \cup P_1)$, $\operatorname{ins}(P \cup P_2)$ includes $\operatorname{ins}(P_1 \cup P_2) = \operatorname{ins}(C)$. This contradicts the fact that $C$ was chosen with $\operatorname{ins}(C)$ maximal, for $\operatorname{ins}(P \cup P_1)$, $\operatorname{ins}(P \cup P_2) \ne \operatorname{ins}(C)$ since they include $U(P)$. We deduce that there is no such $P$, and so (i) holds. In particular, if $\Gamma'$ is non-separable then $U(\Gamma') \subseteq \operatorname{ins}(C)$. By the first slope axiom, $\operatorname{ins}(\Gamma') \subseteq \operatorname{ins}(C)$. Since $\operatorname{ins}(C) \subseteq \operatorname{ins}(\Gamma')$, this proves (iii).

To complete the proof of (ii), let $e \in E(C)$, and let $r$ be the region of $\Gamma$ incident with $e$ with $r \not\subseteq \operatorname{ins}(C)$. We suppose, for a contradiction, that $\Gamma'$ captures $r$. Let $C'$ be a circuit of $\Gamma'$ with $r \subseteq \operatorname{ins}(C')$. If $U(C) \subseteq \operatorname{ins}(C')$, then by the first slope axiom, $\operatorname{ins}(C) \subseteq \operatorname{ins}(C')$, contrary to our choice of $C$. Thus $U(C) \not\subseteq \operatorname{ins}(C')$, and so there are distinct vertices $a, b \in V(C \cap C')$. But by (i), $U(P) \subseteq \operatorname{ins}(C)$ for both paths $P$ of $C'$ with ends $a, b$, and so $U(C) \subseteq \operatorname{ins}(C')$, a contradiction. We deduce that $\Gamma'$ does not capture $r$, as claimed. $\square$

The next result is a strengthening of (4.3).

\stmt
(4.4) \textit{If $\Gamma'$ is confined then $\Sigma - \operatorname{ins}(\Gamma')$ is a region of $\Gamma'$.}

\textit{Proof.} We claim first of all that there is at least one region of $\Gamma'$ not captured by $\Gamma'$. For if $\Gamma'$ has no circuits, then it captures no region, and our claim is true. if $\Gamma'$ has a circuit, we may choose a circuit $C$ with $\operatorname{ins}(C)$ maximal, and our claim follows from (4.3)(ii).

It remains to show that there is at most one such region. We prove by induction on $|E(\Gamma')|$ that if $r'_1, r'_2$ are distinct regions of $\Gamma'$ then $\Gamma'$ captures at least one of them. For if $\Gamma'$ has no circuits, it has only one region (since $\Sigma$ is connected) and our claim is vacuous. Thus we may assume that $\Gamma'$ has a circuit $C$. Let $r'$ be a region of $\Gamma'$ with $r' \subseteq \operatorname{ins}(C)$, and let $e$ be an edge of $\Gamma'$ incident with $r'$. Let $\Gamma''$ be the subdrawing of $\Gamma'$ obtained by deleting $e$. Let $r''_1, r''_2$ be regions of $\Gamma''$ with $r'_1 \subseteq r''_1$, $r'_2 \subseteq r''_2$. Now $\Gamma'$ captures $r'$, and so if $r'_1 = r'$ or $r'_2 = r'$ the claim is true; we assume $r'_1, r'_2 \ne r'$. Since $r'$ is incident with $e$, at least one of $r'_1, r'_2$ is not incident with $e$, say $r'_1$. Hence $r''_1 = r'_1$; but $r'_2 \not\subseteq r'_1$, and so $r''_1 \ne r''_2$. From our inductive hypothesis we deduce that $\Gamma''$ and hence $\Gamma'$ captures at least one of $r''_1, r''_2$, and hence captures at least one of $r'_1, r'_2$, as required. $\square$

\stmt
(4.5) \textit{Let $\Gamma' \subseteq \Gamma$ be confined, and let $r' \in R(\Gamma')$. Suppose that $\Gamma$ captures $r'$. Then so does $\Gamma'$.}

\textit{Proof.} Let $C$ be a circuit of $\Gamma$ with $r' \subseteq \operatorname{ins}(C)$. Then $\operatorname{ins}(C)$ is a closed disc, and there is a circuit $C'$ of $\Gamma'$ bounding a closed disc $\Delta' \subseteq \operatorname{ins}(C)$, with
